\section{First construction: Research U211 to U210}\label{sec:u210}

\subsection{Move the continuation while preserving the Research selector}
Starting with the Research table, apply these eight left associations:
\[
\begin{array}{rrrr}
(162,203,0),&(203,223,0),&(223,188,0),&(188,177,0),\\
(177,89,0),&(89,174,0),&(174,217,0),&(217,163,0).
\end{array}
\]
Here \((p,q,0)\) changes \(((A,B),C)\) at decoder \(p\), whose zero child is
\(q\), into \((A,(B,C))\). The builder verifies pure decoder rows, the complete
unique incoming edge to the inner state, and closure of each of the three
source factors. The independent checker recovers the affected occurrences
by literal source-prefix execution and checks both complete control graphs
after every source change.

These associations carry continuation 116 to the right of state 163.
The Research selector rows at 220 and 129 remain intact.
The old 181 state remains deleted. The final remaining entrances to 218 are
exactly
\[
 158/0:\row0R{218},\qquad
 163/0:\row0R{218},\qquad
 211/0:\row0R{218}.
 \tag{1}
\]
Consequently each visit to 218 from state 0 has a zero immediately left of
the head. This is a direct consequence of the preceding write/move, valid
for every initial tape, rather than a property guessed from a blank trace.

\subsection{The 218-to-157 local identity}
The old rows are \(218/0=\row0R{90}\) and \(218/1=\row1R{86}\).
Redirect all targets of 218 to 157 and delete 218. With the original head at
coordinate zero, the existing table executes
\[
\begin{array}{c|l}
\text{current bit}&\text{replacement state/head path}\\
\hline
0&157@0\to67@{-1}\to151@0\to90@1\\
1&157@0\to148@{-1}\to186@0\to86@1.
\end{array}
\tag{2}
\]
Only cells \(-1,0\) are read. The predecessor is zero by (1);
each write preserves its bit, and the exterior is untouched.

\begin{lemma}[U210 final local fold]
From every common initial tape in state 0, the 211-state prefold table and
the delivered U210 halt equivalently. Each old 218 step becomes three new
steps; every other old step becomes one new step.
\end{lemma}
\begin{proof}
Map 218 to 157 and fix all other labels. For ordinary old rows, write and move
are unchanged and the target is mapped by this projection. For old state 218,
use the exhaustive two cases (2) and the incoming-edge invariant (1).
Every replacement block has positive length and no internal halt.
Induction on old steps gives the simulation, including equality of the whole
tape at its block boundaries. Conversely, if the new machine halts after
\(t\) steps, simulating \(t\) old steps would pass that point unless the old
machine had already halted; positivity and the absence of internal halt
exclude a spurious new halt.
\end{proof}

For any halting old run, or aligned finite run prefixes, this final fold gives
\(T_{\rm new}=T_{\rm old}+2k\), where \(k\) counts old visits to 218, and
\(T_{\rm old}\le T_{\rm new}\le3T_{\rm old}\).
This clock bound is not asserted across the preceding address changes.

The exact U210 dense hash is
\begin{center}\small
\hash{9020e082ef6995d13cd571c325f4d83bf615f4d0a82e77c580353c1127a0d596}.
\end{center}
Its full final-table runtime and blank startup were checked independently of
the state-count calculation.

\section{Second construction: a compatible U209}\label{sec:u209}

\subsection{The shared helper that prevents a direct patch transplant}
In U210, state 129 is entered both by a right-moving source decoder and by
the backward call \(220/0=\row0L{129}\). Its two incoming rows are
\(215/1=\row1R{129}\) and \(220/0=\row0L{129}\).
Treating it as an ordinary singly entered source decoder would lose one
of these uses.

The runtime call begins after 220 has just read and preserved zero.
If the predecessor bit is \(t\), its path is
\[
 220@0\to129@{-1}
 \to\begin{cases}158@0,&t=0,\\204@0,&t=1,\end{cases}
 \to\begin{cases}157@1,&t=0,\\147@1,&t=1.\end{cases}
 \tag{3}
\]
Thus the call uses 129's two rows, but only the zero rows of 158 and 204.
The row \(204/1\) is available for a source-layout change without altering
this runtime block. This selective preservation is broader than requiring
the entire shared decoder to be frozen or to have a single entrance.

\subsection{A closed eight-factor reassociation}
At the U210 source prefix
\[
 \texttt{0011001110101011111111}
\]
state 174 represents primitive ordinals \([3016,3736)\), a region of length 720.
Write its eight closed factors in their unchanged order as follows.

\begin{table}[ht]
\centering\small
\begin{tabular}{cl r ll}
\toprule
Factor & Physical entry & Length & Old suffix & New suffix\\
\midrule
\(A\)&98&26&00&00\\
\(B\)&158&134&010&010\\
\(C\)&147&52&0110&0110\\
\(D\)&146&108&01110&0111\\
\(E\)&220, zero tag&136&011110&10\\
\(F\)&90&27&011111&1100\\
\(G\)&211&103&10&1101\\
\(H\)&163&134&11&111\\
\bottomrule
\end{tabular}
\caption{The zero tag in factor \(E\) is part of its physical decoding context.}
\end{table}

Replace the source shape by
\begin{align*}
 &((A,(B,(C,(D,(E,F))))),(G,H))\\
 &\hspace{2cm}\longmapsto
 ((A,(B,(C,D))),(E,((F,G),H))).
\end{align*}
Every factor is closed and nonhalting. The word and both complete ordinal
control graphs are preserved. Physically, this is implemented by seven row
changes:
\[
\begin{array}{c|c|c}
\text{row}&\text{before}&\text{after}\\
\hline
204/1&\row1R{214}&\row1R{146}\\
214/0&\row0R{146}&\row0R{216}\\
214/1&\row1R{216}&\row1R{163}\\
216/0&\row0R{220}&\row0R{90}\\
216/1&\row1R{90}&\row1R{211}\\
217/0&\row0R{211}&\row0R{220}\\
217/1&\row1R{163}&\row1R{214}
\end{array}
\tag{4}
\]
No state is added. State 216 now decodes the pair \((90,211)\).

The zero-tag entrance to 220 moves from \(216/0\) to \(217/0\); it still
writes zero and moves right. Its one-tag entrance \(185/1\) is unchanged.
All rows used by (3) are unchanged. In particular \(204/0\) is preserved
while \(204/1\) changes. The source entry and backward runtime entry are
therefore justified separately. Exhaustive final-prefix and primitive-runtime
checks additionally validate the actual literal table.

\subsection{Retag state 180 and repair the arithmetic entrance}
Apply the three right associations
\[
 (206,199,1),\qquad(205,206,1),\qquad(7,205,1),
 \tag{5}
\]
where side 1 changes \((A,(B,C))\) to \(((A,B),C)\).
The inner decoder has its specified unique incoming edge in each case, and
the three factors are closed. Both complete ordinal control graphs remain
unchanged.

State 7 also has the backward arithmetic entrance \(171/0=\row0L7\).
Repair this one row to
\[
                     171/0=\row0L{205}.
 \tag{6}
\]
For a current zero at coordinate zero, split on the predecessor:
\[
\begin{array}{c|l|l}
\text{left bit}&\text{old path}&\text{repaired path}\\
\hline
0&171@0\to7@{-1}\to86@0&
 171@0\to205@{-1}\to86@0\\
1&171@0\to7@{-1}\to205@0\to164@1&
 171@0\to205@{-1}\to206@0\to164@1.
\end{array}
\tag{7}
\]
Both paths preserve the two cells and read no exterior cell. This establishes
the entry repair on every such configuration.

The first two stages of (5) have maximum zero run seven; their intermediate
tables are not asserted to be complete runtimes. The final stage restores
maximum zero run six. The independent receipt records any failed intermediate
prefix check as \texttt{NOT\_A\_RUNTIME\_INTERMEDIATE}; it checks every prefix
of the final table successfully.

\subsection{Delete 180 using a five-step contextual path}
In the 210-state prefold table \(R\),
\[
\begin{array}{c|cc}
q&\delta(q,0)&\delta(q,1)\\
\hline
216&\row0R{90}&\row1R{211}\\
180&\row0R{152}&\row1R{211}.
\end{array}
\]
Its complete incoming inventories are
\[
 \operatorname{In}(216)=\{214/0:\row0R{216}\},\qquad
 \operatorname{In}(180)=
 \{132/1:\row1R{180},\,199/1:\row1R{180}\}.
 \tag{8}
\]
Neither state is initial. Therefore the old state is encoded by the bit
immediately left of the head: zero for 216 and one for 180.

Let \(\rho(180)=216\) and \(\rho(q)=q\) otherwise. Redirect both incoming targets
of 180 to 216, delete its two rows, and replace
\[
                           216/0\quad\text{by}\quad\row0L{211}.
 \tag{9}
\]
Keep \(216/1=\row1R{211}\). The resulting machine \(C\) has 209 states.

\begin{table}[ht]
\centering\small
\begin{tabularx}{\linewidth}{@{}rrrrX@{}}
\toprule
Old state & Left bit & Read & New steps & New state/head path\\
\midrule
216&0&0&5&\(216@0,211@{-1},157@0,67@{-1},151@0,90@1\)\\
180&1&0&3&\(216@0,211@{-1},184@0,152@1\)\\
216&0&1&1&\(216@0,211@1\)\\
180&1&1&1&\(216@0,211@1\)\\
\bottomrule
\end{tabularx}
\caption{Exhaustive whole-tape identities for the final fold. Each old block
has one step. Only the predecessor and current cell are read.}
\label{tab:fold}
\end{table}

The five-step case is the compatibility point: 218 has already been deleted,
so the zero predecessor branch must pass through its existing 157 detour.
Using the earlier three-step formula here would be wrong.

\begin{lemma}[Final fold, arbitrary exterior]\label{lem:209fold}
From state 0 on every common initial tape, \(R\) and \(C=U_{209}\) halt
equivalently. For a halting old run,
\[
 T_C=T_R+4n_0+2n_1,\qquad T_R\le T_C\le5T_R,
\]
where \(n_0\) counts old read-zero visits to 216 and \(n_1\) counts old
read-zero visits to 180.
\end{lemma}
\begin{proof}
Inventory (8) establishes the tag invariant whenever either old state is
entered. In the four exhaustive cases of Table~\ref{tab:fold}, direct literal
execution leaves both displayed cells unchanged, reads no other cell, and
ends with the same head and the projected old target. No internal state is
the halt label. Every remaining old row commutes with \(\rho\), preserving
its write and direction. Induction gives a positive block simulation.
A new finite halt cannot occur inside a block, and each block consumes at
least one new step, so the converse follows by the finite-horizon argument
used for U210. Summing the four block lengths gives the clock identity.
For arithmetic codes of finite-support tapes these are finite induction
arguments formalizable in PA.
\end{proof}

The all-initial-tape assertion and time bound concern \(R\to C\) alone.
They are not extended to the preceding source-address transformations.

Relative to U210, U209 changes 15 retained rows and deletes both rows of 180.
It changes 847 source addresses, including 508 changed by (4), while retaining
the same 197 distinct syntactic subtrees.
